% Formato 3 — notas no corpo (160×230 mm).  O texto oficial vem em destaque, em corpo
% maior (12 pt, preto); cada nota é solta logo DEPOIS do parágrafo que a chama (não no
% meio da frase), em corpo menor, cinza, recuada e com filete à esquerda.
\documentclass[12pt,twoside]{book}
\usepackage[paperwidth=160mm,paperheight=230mm,inner=19mm,outer=24mm,top=17mm,
            bottom=19mm,includehead,headsep=6mm,footskip=9mm,marginparwidth=14mm]{geometry}
% Macros comuns aos três ensaios de diagramação (formato-*.tex).  Cada mestre define a
% geometria e o tratamento das notas ANTES de % Macros comuns aos três ensaios de diagramação (formato-*.tex).  Cada mestre define a
% geometria e o tratamento das notas ANTES de % Macros comuns aos três ensaios de diagramação (formato-*.tex).  Cada mestre define a
% geometria e o tratamento das notas ANTES de \input{comum}; aqui ficam fonte, títulos,
% marcas de página e cabeçalhos.
\usepackage{fontspec}
\usepackage{polyglossia}
\setmainlanguage{portuguese}
\setotherlanguage{latin}
\setmainfont{EB Garamond}[Numbers={OldStyle,Proportional}]
\usepackage{microtype}
\usepackage{xcolor}
\definecolor{marca}{gray}{0.45}
\setlength{\emergencystretch}{1.5em}
\frenchspacing

% --- página do original ------------------------------------------------------
% \pg{443}: virada de página do original no corpo -> na margem externa.
% \pgn{444}: virada dentro de nota (marginpar é proibido ali) -> discreta no fluxo.
% Em modo vertical (antes de um título ou parágrafo) a marca espera o próximo
% parágrafo, senão cairia no pé da página anterior.
\newcommand{\margempg}[1]{\marginpar{\footnotesize\color{marca}[#1]}}
\gdef\pgpendente{}
\newcommand{\pg}[1]{%
  \ifvmode\xdef\pgpendente{\unexpanded\expandafter{\pgpendente}\noexpand\margempg{#1}}%
  \else\margempg{#1}\fi\ignorespaces}
\AddToHook{para/begin}{\ifinner\else
  \let\pgagora\pgpendente\gdef\pgpendente{}\pgagora\fi}
\newcommand{\pgn}[1]{{\scriptsize\color{marca}[#1]}}
% (classe book, twoside: \marginpar vai para a margem externa)
\setlength{\marginparsep}{3mm}

% --- estrutura do autor ------------------------------------------------------
% \capitulo{rótulo}{título}{summario}{pág.}
\newcommand{\capitulo}[4]{%
  \clearpage\thispagestyle{plain}\pg{#4}%
  \vspace*{2\baselineskip}%
  \begin{center}
    {\scshape\lowercase{#1}\par}\medskip
    {\large\scshape\lowercase{#2}\par}
  \end{center}
  \medskip
  {\small\noindent\textsc{summario}: --- #3\par}
  \markboth{\lowercase{#2}}{\lowercase{#2}}%
  \bigskip}
% \artigo{rótulo}{título}
\newcommand{\artigo}[2]{%
  \bigskip
  \begin{center}{\scshape\lowercase{#1}.}\\[.3ex]{\itshape #2}\end{center}
  \medskip}
% tese numerada que corre com o texto / em linha própria / subdivisão a), b)
\newcommand{\tese}[2]{\par\medskip\noindent#1.\enspace\textbf{#2}\;---\;}
\newcommand{\teseB}[2]{\par\medskip\noindent#1.\enspace\textbf{#2}\par}
\newcommand{\subp}[1]{\par\noindent\hspace*{\parindent}#1)\enspace}
% trecho omitido neste ensaio
\newcommand{\omitido}{\par{\centering[\dots]\par}}

% --- cabeçalho ---------------------------------------------------------------
\usepackage{fancyhdr}
\pagestyle{fancy}
\fancyhf{}
\fancyhead[LE]{\small\oldstylenums{\thepage}\quad\scshape theodicea}
\fancyhead[RO]{\small\scshape essencia e attributos de deus\quad\oldstylenums{\thepage}}
\renewcommand{\headrulewidth}{0pt}
\fancypagestyle{plain}{\fancyhf{}\fancyfoot[C]{\small\oldstylenums{\thepage}}}
\setcounter{page}{443}
; aqui ficam fonte, títulos,
% marcas de página e cabeçalhos.
\usepackage{fontspec}
\usepackage{polyglossia}
\setmainlanguage{portuguese}
\setotherlanguage{latin}
\setmainfont{EB Garamond}[Numbers={OldStyle,Proportional}]
\usepackage{microtype}
\usepackage{xcolor}
\definecolor{marca}{gray}{0.45}
\setlength{\emergencystretch}{1.5em}
\frenchspacing

% --- página do original ------------------------------------------------------
% \pg{443}: virada de página do original no corpo -> na margem externa.
% \pgn{444}: virada dentro de nota (marginpar é proibido ali) -> discreta no fluxo.
% Em modo vertical (antes de um título ou parágrafo) a marca espera o próximo
% parágrafo, senão cairia no pé da página anterior.
\newcommand{\margempg}[1]{\marginpar{\footnotesize\color{marca}[#1]}}
\gdef\pgpendente{}
\newcommand{\pg}[1]{%
  \ifvmode\xdef\pgpendente{\unexpanded\expandafter{\pgpendente}\noexpand\margempg{#1}}%
  \else\margempg{#1}\fi\ignorespaces}
\AddToHook{para/begin}{\ifinner\else
  \let\pgagora\pgpendente\gdef\pgpendente{}\pgagora\fi}
\newcommand{\pgn}[1]{{\scriptsize\color{marca}[#1]}}
% (classe book, twoside: \marginpar vai para a margem externa)
\setlength{\marginparsep}{3mm}

% --- estrutura do autor ------------------------------------------------------
% \capitulo{rótulo}{título}{summario}{pág.}
\newcommand{\capitulo}[4]{%
  \clearpage\thispagestyle{plain}\pg{#4}%
  \vspace*{2\baselineskip}%
  \begin{center}
    {\scshape\lowercase{#1}\par}\medskip
    {\large\scshape\lowercase{#2}\par}
  \end{center}
  \medskip
  {\small\noindent\textsc{summario}: --- #3\par}
  \markboth{\lowercase{#2}}{\lowercase{#2}}%
  \bigskip}
% \artigo{rótulo}{título}
\newcommand{\artigo}[2]{%
  \bigskip
  \begin{center}{\scshape\lowercase{#1}.}\\[.3ex]{\itshape #2}\end{center}
  \medskip}
% tese numerada que corre com o texto / em linha própria / subdivisão a), b)
\newcommand{\tese}[2]{\par\medskip\noindent#1.\enspace\textbf{#2}\;---\;}
\newcommand{\teseB}[2]{\par\medskip\noindent#1.\enspace\textbf{#2}\par}
\newcommand{\subp}[1]{\par\noindent\hspace*{\parindent}#1)\enspace}
% trecho omitido neste ensaio
\newcommand{\omitido}{\par{\centering[\dots]\par}}

% --- cabeçalho ---------------------------------------------------------------
\usepackage{fancyhdr}
\pagestyle{fancy}
\fancyhf{}
\fancyhead[LE]{\small\oldstylenums{\thepage}\quad\scshape theodicea}
\fancyhead[RO]{\small\scshape essencia e attributos de deus\quad\oldstylenums{\thepage}}
\renewcommand{\headrulewidth}{0pt}
\fancypagestyle{plain}{\fancyhf{}\fancyfoot[C]{\small\oldstylenums{\thepage}}}
\setcounter{page}{443}
; aqui ficam fonte, títulos,
% marcas de página e cabeçalhos.
\usepackage{fontspec}
\usepackage{polyglossia}
\setmainlanguage{portuguese}
\setotherlanguage{latin}
\setmainfont{EB Garamond}[Numbers={OldStyle,Proportional}]
\usepackage{microtype}
\usepackage{xcolor}
\definecolor{marca}{gray}{0.45}
\setlength{\emergencystretch}{1.5em}
\frenchspacing

% --- página do original ------------------------------------------------------
% \pg{443}: virada de página do original no corpo -> na margem externa.
% \pgn{444}: virada dentro de nota (marginpar é proibido ali) -> discreta no fluxo.
% Em modo vertical (antes de um título ou parágrafo) a marca espera o próximo
% parágrafo, senão cairia no pé da página anterior.
\newcommand{\margempg}[1]{\marginpar{\footnotesize\color{marca}[#1]}}
\gdef\pgpendente{}
\newcommand{\pg}[1]{%
  \ifvmode\xdef\pgpendente{\unexpanded\expandafter{\pgpendente}\noexpand\margempg{#1}}%
  \else\margempg{#1}\fi\ignorespaces}
\AddToHook{para/begin}{\ifinner\else
  \let\pgagora\pgpendente\gdef\pgpendente{}\pgagora\fi}
\newcommand{\pgn}[1]{{\scriptsize\color{marca}[#1]}}
% (classe book, twoside: \marginpar vai para a margem externa)
\setlength{\marginparsep}{3mm}

% --- estrutura do autor ------------------------------------------------------
% \capitulo{rótulo}{título}{summario}{pág.}
\newcommand{\capitulo}[4]{%
  \clearpage\thispagestyle{plain}\pg{#4}%
  \vspace*{2\baselineskip}%
  \begin{center}
    {\scshape\lowercase{#1}\par}\medskip
    {\large\scshape\lowercase{#2}\par}
  \end{center}
  \medskip
  {\small\noindent\textsc{summario}: --- #3\par}
  \markboth{\lowercase{#2}}{\lowercase{#2}}%
  \bigskip}
% \artigo{rótulo}{título}
\newcommand{\artigo}[2]{%
  \bigskip
  \begin{center}{\scshape\lowercase{#1}.}\\[.3ex]{\itshape #2}\end{center}
  \medskip}
% tese numerada que corre com o texto / em linha própria / subdivisão a), b)
\newcommand{\tese}[2]{\par\medskip\noindent#1.\enspace\textbf{#2}\;---\;}
\newcommand{\teseB}[2]{\par\medskip\noindent#1.\enspace\textbf{#2}\par}
\newcommand{\subp}[1]{\par\noindent\hspace*{\parindent}#1)\enspace}
% trecho omitido neste ensaio
\newcommand{\omitido}{\par{\centering[\dots]\par}}

% --- cabeçalho ---------------------------------------------------------------
\usepackage{fancyhdr}
\pagestyle{fancy}
\fancyhf{}
\fancyhead[LE]{\small\oldstylenums{\thepage}\quad\scshape theodicea}
\fancyhead[RO]{\small\scshape essencia e attributos de deus\quad\oldstylenums{\thepage}}
\renewcommand{\headrulewidth}{0pt}
\fancypagestyle{plain}{\fancyhf{}\fancyfoot[C]{\small\oldstylenums{\thepage}}}
\setcounter{page}{443}

\usepackage[most]{tcolorbox}
\definecolor{nota}{gray}{0.28}
\definecolor{filete}{gray}{0.70}

% \footnote{...}: chamada «(n)» no texto; o conteúdo entra numa fila esvaziada no
% fim do parágrafo (hook para/after).  \ifemnota impede que os \par da própria nota
% disparem de novo a fila.
\newtoks\notasfila
\newif\ifemnota
\renewcommand\footnote[1]{%
  \stepcounter{footnote}%
  \hbox{\normalfont\,(\arabic{footnote})}%
  \expandafter\enfileira\expandafter{\the\value{footnote}}{#1}}
\newcommand\enfileira[2]{\global\notasfila\expandafter{\the\notasfila\umanota{#1}{#2}}}
\newcommand\umanota[2]{%
  \begin{tcolorbox}[breakable,enhanced jigsaw,frame hidden,interior hidden,
      borderline west={1.2pt}{0pt}{filete},left=4mm,right=0mm,top=0mm,bottom=0mm,
      boxsep=0mm,before skip=.6\baselineskip,after skip=.9\baselineskip,
      left skip=6mm,fontupper=\footnotesize\color{nota}]
    \setlength{\parindent}{1em}\noindent{\scshape nota (#1).}\enspace #2
  \end{tcolorbox}}
% Só na lista vertical principal: cabeçalho, marginpar e a própria caixa são
% "inner" e também disparam para/after (sem isto a fila era solta dentro da rotina de
% saída quando o parágrafo da chamada quebrava de página, e a nota saía sobreposta).
\newcommand\soltanotas{%
  \ifinner\else\ifemnota\else
    \edef\tmpfila{\the\notasfila}%
    \ifx\tmpfila\empty\else
      \emnotatrue
      \begingroup\the\notasfila\endgroup
      \global\notasfila{}%
      \emnotafalse
    \fi
  \fi\fi}
\AddToHook{para/after}{\soltanotas}
\begin{document}
% Trecho de exemplo para os ensaios de diagramação: vol. 2, Theodicea, cap. II,
% Art. I (nn. 22–25), pp. 443–449 do original.
% Base: revisao/essencia-de-deus/revisado.tex (TEI revisado por LLM), com os erros de
% OCR evidentes corrigidos à mão (perfeicoes, eztrinseca, Dus, alieujus, ouires, c. 48→43…).
% Ortografia de 1927 mantida.  Ainda NÃO conferido linha a linha contra a imagem.
% Omitido (marcado […]): o fim da nota de 24 c), «Os erros ácerca da infinidade…»,
% cuja fronteira corpo/nota ficou em dúvida na revisão (relatorio.md, p. 445).
%
% Macros esperadas do documento-mestre: \pg, \pgn, \capitulo, \artigo, \tese, \sub.

\capitulo{Capitulo segundo}{Essencia e attributos de Deus}%
  {Essencia de \textsc{Deus}, seu constitutivo e attributos. --- Attributos de
   \textsc{Deus} negativos. --- Attributos de \textsc{Deus} positivos.}{443}

\artigo{Artigo I}{Essencia de Deus, seu constitutivo e attributos.}

\tese{22}{Essencia de Deus.} \emph{Essencia} de \textsc{Deus} é o que faz com que Elle
seja \textsc{Deus}, e se distinga de todos os outros seres.

\tese{23}{Essencia physica e metaphysica de Deus.} A essencia \emph{physica} de
\textsc{Deus} é constituida por todas as perfeições, de que a Divindade é dotada. A
\emph{metaphysica} é a perfeição fundamental, que \emph{se concebe} como raiz e fonte de
todas as outras perfeições divinas, e pela qual \textsc{Deus} se distingue das
creaturas. --- A essencia \emph{physica} representa \textsc{Deus} como é em si mesmo; a
\emph{metaphysica} exprime \textsc{Deus} como é em relação ao nosso modo de
entender\footnote{Lembramos que a essencia --- é \emph{physica}, quando os seus elementos
constitutivos são \emph{reaes} e \emph{concretos} (assim o \emph{corpo organico} e a
\emph{alma racional} constituem a essencia \emph{physica} do homem), --- é
\emph{metaphysica}, quando os seus elementos constitutivos são \emph{logicos} e
\emph{abstractos}, denotando o \emph{genero} e a \emph{differença} (assim a
\emph{animalidade} e a \emph{racionalidade} constituem a essencia \emph{metaphysica} do
homem). --- Da essencia brotam as \emph{propriedades}. Estas distinguem-se
\emph{realmente} da essencia \emph{physica}, porque não a constituem mas resultam d'ella
(assim a intelligencia e a vontade são realmente distinctas da essencia \emph{physica} da
nossa alma), --- mas podem não distinguir-se \emph{realmente} da essencia
\emph{metaphysica}, porque, ainda que não possam considerar-se como elementos
\emph{constitutivos} da essencia, comtudo podem identificar-se com esses elementos (assim
a \emph{immortalidade}, embora não possa considerar-se como um elemento constitutivo da
essencia \emph{metaphysica} da nossa alma, comtudo não se distingue \emph{realmente} da
mesma essencia). E' claro que não podemos fallar da essencia de \textsc{Deus}, como
fallamos da essencia das creaturas. Em Deus, a essencia \emph{physica} significa e
abrange toda a perfeição, que compete ao Ente infinito, sem sombra de imperfeição, --- e
a essencia \emph{metaphysica} exclue toda e qualquer composição. Por isso o
\emph{constitutivo} da essencia \emph{metaphysica} de Deus não é uma coisa realmente
distincta das outras propriedades, e a sua relação com estas deve, mais ou menos,
conceber-se como se concebe a relação do ente, considerado na sua razão communissima, com
as suas propriedades \emph{transcendentaes}.

Dizemos que a essencia \emph{metaphysica} exprime \textsc{Deus} como Elle é em
\emph{relação} ao \emph{nosso modo de entender}. Por quanto, considerando os seres
creados, \pgn{444}~encontramos em cada um d'elles uma perfeição, que é a primeira a ser
conhecida, que constitue o mesmo ser numa determinada especie, distinguindo-o de todos os
outros, e que é a raiz e a fonte de todas as outras perfeições. E, como subimos ao
conhecimento de \textsc{Deus} por meio das creaturas, procuramos encontrar em
\textsc{Deus} essa perfeição \emph{primeira}, \emph{caracteristica e radical}; embora na
essencia de \textsc{Deus}, como é em si mesma, não se encontre attributo algum, que seja
anterior aos outros attributos e de que estes fluam. --- Este modo de entender a essencia
de \textsc{Deus} é o resultado da limitação da nossa intelligencia; não podemos entender o
infinito senão por partes, e, ainda assim, com muita imperfeição. Todavia, a propria
intelligencia emenda ou explica o seu modo de ver; pois que, depois de ter considerado um
attributo divino, como fundamento dos outros, prova que, em \textsc{Deus}, não só não ha
distincção alguma entre a essencia e os attributos, mas nem sequer entre os mesmos
attributos.}.

\pg{444}\teseB{24}{A essencia physica de Deus consiste na infinidade.}

\subp{a} Em Deus, o ser (ou a existencia) não é recebido nem subsiste na essencia, mas é
\emph{subsistente por si}; visto que o ser de Deus não se distingue da sua essencia, e o
que não se distingue de uma coisa, não pode ser recebido nem subsistir nella. Ora o ente,
cujo ser não é recebido nem subsiste na essencia, é necessariamente infinito; porque o ser
é limitado ou finito só porque e quando é recebido numa certa e determinada essencia. Logo
Deus é infinito, isto é, a essencia \emph{physica} de \textsc{Deus} consiste na
\emph{infinidade}\footnote{Desenvolvamos este argumento. Uma substancia, como dissemos na
\emph{Ontologia}, pode \emph{subsistir}, ou \emph{existir em si}, por dois modos: ou
porque é a existencia, ou porque \emph{tem} a existencia. No primeiro caso, a existencia
(ou o ser) não é recebida na essencia, ou natureza, como na sua potencia proporcionada, e
por isso é acto puro e subsistente por si; --- no segundo, a existencia (ou o ser) é
recebida na essencia, como acto e principio de subsistencia, e por isso não é acto puro,
nem subsiste por si, mas subsiste na essencia, como na sua potencia. Ora, em Deus, a
essencia e o ser, ou a existencia, não se distinguem \emph{realmente}, porque o ser não é
accidental, mas é essencial á essencia divina; porque esta é \emph{essencialmente}
existente. Se o ser não se distingue \emph{realmente} da essencia divina, Deus, assim como
é a sua essencia, isto é, a sua Divindade, assim tambem é o seu proprio ser, ou a sua
propria existencia. Mas o ser, que não se distingue da essencia, não é recebido na
propria essencia; porque o ser, que é recebido na essencia, deve necessariamente
distinguir-se d'esta, como o acto se distingue necessariamente da potencia (assim a côr,
recebida numa parede, distingue-se da propria parede). O ser, que não é recebido numa
essencia, não é limitado ou determinado a uma certa especie; visto que o ser emtanto se
limita ou se determina a uma certa especie, emquanto é recebido numa certa potencia ou
natureza (assim o ser \emph{humano} é limitado ou determinado a uma certa especie, e não a
uma outra, porque é recebido na \emph{essencia} ou \emph{natureza humana}). Portanto, o
ser, que não é recebido numa essencia, não é limitado ou circumscripto, mas é illimitado
ou infinito. E, como o \emph{ser}, por ser \emph{acto}, \pgn{445}~indica e é
\emph{perfeição}, segue-se que o ser subsistente por si é a perfeição subsistente,
illimitada, infinita; de modo que não ha nem pode conceber-se entitade, perfeição,
attributo, que esse ser não possua, assim como a alvura, se fosse subsistente em si, havia
de conter tudo o que tem razão de alvor. Deus, portanto, sendo o ser subsistente por si,
é infinito na sua perfeição. E' o raciocinio do Doutor Angelico, quando diz: «~Deus est
ipsum esse per se subsistens; ex quo oportet quod totam perfectionem essendi in se
contineat. Manifestum est enim quod si aliquod calidum non habeat totam perfectionem
calidi, hoc ideo est, quia calor non participatur secundum perfectam rationem; sed, si
calor esset per se subsistens, non posset ei aliquid deesse de virtute caloris. Unde cum
Deus sit ipsum esse subsistens, nihil de perfectione essendi potest ei deesse. Omnium
autem perfectiones pertinent ad perfectionem essendi; secundum hoc enim aliqua perfecta
sunt quod aliquo modo esse habeant. Unde sequitur quod nullius rei perfectio Deo desit~»
(\emph{Sum. Theolog.}, p.~I, q.~4, a.~2). --- Mais adiante accrescenta: «~Cum esse
divinum non sit esse receptum in aliquo; sed ipse sit suum esse subsistens, manifestum est
quod ipse Deus sit infinitus et perfectus~» (\emph{Ib.}, q.~7, a.~1). E noutro logar:
«~Deus est infinitus: quod sic videri potest. Esse enim hominis terminatum est ad hominis
speciem, quia est receptum in natura speciei humanae, et simile est de esse equi, vel
cujuslibet creaturae. Esse autem Dei, quum non sit in aliquo receptum, sed sit esse
purum, non limitatur ad aliquem modum perfectionis essendi, sed totum esse in se habet~»
(\emph{De Pot.}, q.~1, a.~2).}.

\pg{445}\subp{b} A essencia \emph{physica} de \textsc{Deus} é constituida por todas as
perfeições, de que a Divindade é dotada. Ora as perfeições, de que a Divindade é dotada,
são \emph{infinitas}; visto que \textsc{Deus}, sendo \emph{immutavel} e
\emph{necessario}, deve possuir todas as possiveis perfeições; pois, se Lhe faltasse
alguma, seria, com relação a esta, \emph{mutavel} e \emph{contingente}, e por isso
dependeria de uma causa. Logo a essencia \emph{physica} de \textsc{Deus} consiste na
\emph{infinidade}.

\subp{c} Se as perfeições, que constituem a essencia \emph{physica} de \textsc{Deus},
fossem \emph{finitas} ou \emph{limitadas}, essa limitação deveria provir --- ou de uma
causa \emph{extrinseca}, --- ou da indole da propria essencia divina, em que as perfeições
são recebidas. Ora não pode provir --- nem de uma causa \emph{extrinseca}, porque Deus,
sendo a Causa \emph{primeira}, não depende de outra causa, --- nem da indole da propria
essencia divina, porque esta, sendo perfeita por \emph{si mesma}, e não por
\emph{participação}, não pode deixar de possuir a plenitude da perfeição. Logo as
perfeições, que constituem a essencia \emph{physica} de \textsc{Deus}, são
\emph{infinitas}. Logo a essencia \emph{physica} de \textsc{Deus} consiste na
\emph{infinidade}\footnote{A \emph{infinidade} de \textsc{Deus} \emph{indirectamente}
demonstra-se pelo facto de se não encontrar n'Elle nenhuma das causas, que limitam as
creaturas. Na verdade, as creaturas, que são compostas de \emph{materia} e \emph{forma},
são limitadas, porque a \emph{forma} limita a \emph{materia}, quando a determina a
constituir uma certa especie e não outra, --- e a \emph{materia} limita a \emph{forma},
tornando-a \emph{sin\-}\pgn{446}\emph{gular} e \emph{individua}. Ora \textsc{Deus} não
pode ser limitado d'este modo: porque, como provaremos d'aqui a pouco, não é composto da
\emph{materia} e \emph{forma}. --- As creaturas todas pertencem a um certo \emph{genero} e
a uma certa \emph{especie}, e por isso são limitadas, porque não possuem as perfeições
proprias dos outros generos e das outras especies. Ora \textsc{Deus}, sendo o ser
\emph{necessario}, \emph{immutavel} e \emph{independente}, transcende todos os generos e
todas as especies. --- As creaturas todas são compostas de \emph{essencia} e de
\emph{ser}, ou \emph{existencia}, e por isso são limitadas, porque o \emph{ser}, sendo
recebido numa essencia \emph{finita}, como é a essencia da creatura, não pode deixar de
ser \emph{finito}. Ora, como dissemos, o ser divino, \emph{existindo} por assim o exigir a
sua \emph{essencia}, subsiste por \emph{si} e não é recebido numa essencia, que o limite.
--- As creaturas são limitadas pelo \emph{espaço} e pelo \emph{tempo}. \textsc{Deus}, por
ser \emph{immenso}, não é limitado pelo \emph{espaço}, e, por ser \emph{eterno}, não é
determinado pelo \emph{tempo}. Logo \textsc{Deus} é \emph{infinito}. --- S.~Thomaz prova a
mesma verdade, dizendo que Deus é \emph{infinito}; porque a razão da limitação num ente
deriva ou da \emph{causa} do proprio ente ou da \emph{essencia}, em que o ser é recebido,
mas Deus não tem causa e o seu ser não é recebido na essencia, mas, como provámos, é
subsistente por si. Eis as suas palavras: «~Ipsum esse, absolute consideratum, infinitum
est\dots{} Si igitur alicujus esse sit finitum, oportet quod limitetur esse illud per
aliquid aliud, quod sit aliqualiter causa illius esse, vel \emph{receptivum} ejus. Sed
esse divini non potest esse aliqua causa, quia ipse necesse est esse per seipsum: nec esse
ejus est receptivum, quum ipse sit suum esse. Ergo esse suum est infinitum, et ipse
infinitus~» (\emph{C. Gent.}, lib.~I, c.~43).

\omitido}.

\tese{25}{A essencia metaphysica de Deus consiste na subsistencia do proprio ser.} A
essencia \emph{metaphysica} de \textsc{Deus} consiste naquella perfeição, --- que se
concebe como raiz e fonte de todas as outras perfeições divinas, e pela qual \textsc{Deus}
se distingue das creaturas. Ora a \emph{subsistencia do proprio ser} é a
perfeição\pg{446}\pg{447} --- que se concebe como raiz e fonte de todas as outras
perfeições divinas, porque, como provámos, Deus possue todas as perfeições, possiveis e
imaginaveis, pelo facto de ser o \emph{proprio ser subsistente}, --- e pela qual se
distingue das creaturas, visto que Deus é o Ente infinito, porque \emph{é} o proprio ser
subsistente, e as creaturas são finitas, porque \emph{têm} o ser recebido numa
determinada essencia. Logo a essencia \emph{metaphysica} de \textsc{Deus} consiste na
\emph{subsistencia do proprio ser}\footnote{A conclusão é certa, e é uma consequencia dos
principios expostos no numero anterior. O ser divino, pelo facto de não ser recebido em
nenhuma \pgn{448}~essencia ou potencia, não tem limites, encerra uma perfeição absoluta sem
sombra de imperfeição. Repetimos com S.~Thomaz: «~\emph{Ex hoc quod Deus est ipsum esse
subsistens, oportet quod totam perfectionem essendi in se contineat}~» (\emph{Sum. Th.},
p.~I, q.~4, a.~2). Sendo subsistente, o ser divino distingue-se das creaturas, que não
\emph{são} o ser, mas \emph{têm} o ser, recebido numa determinada potencia ou essencia e
por isso limitado a uma certa especie. Como o ser \emph{é acto}, o ser que não é recebido
numa potencia, mas que é acto \emph{puro}, é necessariamente infinito ou illimitado; ao
passo que o ser, que é recebido numa potencia, é limitado por ella. Por isso, não só pode
dizer-se que o \emph{constitutivo} da essencia divina consiste em que \textsc{Deus}
\emph{é o ser subsistente por si}, mas tambem pode dizer-se que esse \emph{constitutivo}
consiste em que \textsc{Deus} é \emph{acto purissimo}, \emph{não recebido em nenhuma
potencia}. As expressões são synonimas. --- Concluimos --- que a razão formal e
intrinseca, pela qual \textsc{Deus} é unico, eterno, infinito, é porque Elle é o proprio
ser subsistente, é acto purissimo, cuja essencia é constituida pelo proprio ser, --- e
que, pelo contrario, a razão, pela qual as creaturas são multiplicaveis, imperfeitas,
finitas, é porque têm o ser recebido na essencia, e por isso são compostas de potencia e
de acto, de essencia e de ser.

Sendo Deus o proprio ser subsistente, o nome mais proprio de Deus é \textsc{Jehovah}, que
significa: \textsc{Qui est} --- \textsc{Aquelle que é}. Por quanto, os nomes significam a
forma ou essencia das coisas. Ora, em Deus, a essencia identifica-se com o ser, isto é, o
ser é a sua propria essencia, de modo que á pergunta: \emph{o que é Deus}? podemos
responder: \emph{Deus é o proprio ser} --- \emph{Deus est ipsum esse}. Logo o nome
\emph{Qui est} é muito proprio de Deus. Com effeito, de nenhuma creatura podemos dizer que
ella é o proprio ser; porque o \emph{ser}, nas \emph{creaturas}, não é a \emph{essencia},
nem é um \emph{predicado essencial}, mas é distincto da essencia, embora seja recebido
nella; e por isso os nomes, que significam as coisas creadas, exprimem a essencia, e não e
ser, assim \emph{homem} exprime a essencia de um individuo, isto é, a \emph{humanidade},
mas não exprime o \emph{ser} ou a \emph{existencia} d'esse individuo, que poderia não
existir. Escutemos S.~Thomaz: «~Cum in omni quod est, sit considerare \emph{quidditatem
suam}, per quam subsistit in natura determinata, et \emph{esse suum}, per quod dicitur de
eo quod est in actu, hoc nomen \emph{res} imponitur rei a quidditate sua, hoc nomen
\emph{qui est} vel \emph{ens} imponitur ab ipso actu essendi. Cum autem ita sit quod in
qualibet re creata, essentia sua differat a suo esse, res illa proprie denominatur a
quidditate sua et non ab actu essendi, sicut homo ab humanitate. In Deo autem ipsum esse
suum est sua quidditas; et ideo nomen quod sumitur ab esse, proprie nominat ipsum, et est
proprium nomen ejus; sicut proprium nomen hominis quod sumitur a quidditate sua~»
(I~\emph{Sent.}, Dist.~8, q.~1, a.~1).

E não pode deixar de ser assim. Quando digo: \emph{Pedro é homem}, attribuo a esse
individuo duas coisas: o \emph{ser}, ou a \emph{existencia}, e a \emph{natureza humana},
isto é, attribuo-lhe o ser limitado e determinado a uma certa especie; mas quando digo:
\emph{Deus é aquelle que é,} quero affirmar que Deus não é este ou aquelle ente, mas que é
o ser, todo o ser, o proprio ser, o ser infinito, de modo que tudo o que tem razão de ser
deve encontrar-se e se encontra \pgn{449}~em Deus. Eis o sentido do nome \emph{qui est},
que o proprio Deus manifestou a Moyses. Lê-se no \emph{Exodo}: «~Ait Moyses ad Deum: Ecce
ego vadam ad filios Israel, et dicam eis: Deus patrum vestrorum misit me ad vos. Si
dixerint mihi: quod est nomen ejus? quid dicam eis? Dixit Deus ad Moysen: \textsc{Ego sum
qui sum}. Ait: Sic dices filiis Israel: \textsc{Qui est} misit me ad vos~» (Cap.~III,
vv.~13, 14). --- E' este um nome mais universal e menos determinado que os outros; mas por
isso mesmo é o mais proprio. Sendo a essencia divina infinitamente superior á
intelligencia creada, o nome mais proprio de Deus é o que melhor traduz essa
transcendencia e indefinibilidade, e exclue as imperfeitas determinações das creaturas.
Diz S.~Thomaz: «~Quia divina comprehendi a nobis non possunt, convenientius a nobis
significantur per nomina communia, quae indefinite aliquid significant, quam per nomina
specialia, quae definite rei species exprimunt; unde hoc nomen \emph{qui est}, quod
secundum Damascenum significat substantiae pelagus infinitum, convenientissimum nomen
dicitur esse~» (\emph{De pot.}, q.~10, a.~1 ad~9).

Nem se diga que o nome \emph{qui est} não é o mais proprio de Deus, porque não o distingue
das creaturas, pois Deus é e as creaturas tambem \emph{são}. Por quanto, ha uma grande
differença entre o \emph{ser} de Deus e o \emph{ser} das creaturas. O ser de Deus, por ser
subsistente em si, identifica-se com a essencia, e por isso se torna individuo e distincto
de todos os outros entes; ao passo que o ser das creaturas, por não ser subsistente em si,
não se identifica com a essencia, mas é recebido e subsiste na essencia, e por isso se
torna individuo e distincto por causa da mesma essencia. Portanto o nome \emph{qui est} é
exclusivamente proprio de Deus, porque só em Deus o ser é a mesma essencia, e o nome, como
dissemos, exprime a essencia. Ouçamos de novo S.~Thomaz: «~Ipsum esse Dei distinguitur et
individuatur a quolibet alio esse per hoc ipsum, quod est esse per se subsistens, et non
adveniens alicui naturae, quae sit aliud ab ipso esse. Omne autem aliud esse, quod non est
subsistens, oportet quod individuetur per naturam et substantiam, quae in tali esse
subsistit. Et in eis verum est, quod esse hujus est aliud ab esse illius per hoc quod est
alterius naturae~» (\emph{De pot.}, q.~7, a.~2 ad~5).}.

\end{document}
